\section{Conclusiones y trabajo futuro}
A partir de la metodología cualitativa implementada en este proyecto, se evidenció que la combinación de una arquitectura sólida y una metodología ágil resulta fundamental para el éxito en el desarrollo de sistemas. La revisión documental mediante el thesaurus permitió identificar conceptos clave como la Programación Orientada a Objetos (POO), el patrón Modelo–Vista–Controlador (MVC) y la metodología Scrum, los cuales aportaron lineamientos claros y prácticos para la construcción del sistema de gestión de eventos.

Los resultados obtenidos demuestran que la aplicación de MVC y POO facilita la organización, escalabilidad y mantenibilidad del software, mientras que el uso de Scrum fortaleció la comunicación del equipo, la planificación de tareas y la entrega incremental de funcionalidades. En conjunto, estos hallazgos confirmaron que la investigación documental y la práctica de desarrollo estuvieron alineadas, logrando así un proyecto formativo con bases teóricas sólidas y un producto funcional que responde a las necesidades planteadas.